\documentclass[11pt]{amsart}
\usepackage{amsmath,amssymb,amsthm}
\usepackage[margin=1.1in]{geometry}
\usepackage{newtxtext,newtxmath}
\usepackage{microtype}
\usepackage[colorlinks=true, linkcolor=blue, citecolor=blue, urlcolor=blue]{hyperref}
\usepackage{listings}
\usepackage{upquote}
\lstdefinelanguage{Lean}{
  morekeywords={def,theorem,abbrev,open,namespace,end},
  sensitive=true,
  morecomment=[l]{--},
}
\lstset{language=Lean, basicstyle=\small\fontfamily{lmtt}\selectfont\upshape, keywordstyle=\bfseries,
  columns=fullflexible, keepspaces=true, breaklines=false,
  xleftmargin=1em, aboveskip=6pt, belowskip=6pt,
  literate={ℕ}{{$\mathbb{N}$}}1 {ℝ}{{$\mathbb{R}$}}1 {ℤ}{{$\mathbb{Z}$}}1
    {→}{{$\rightarrow$}}1 {×}{{$\times$}}1 {∃}{{$\exists$}}1 {∀}{{$\forall$}}1
    {κ}{{$\kappa$}}1 {≤}{{$\le$}}1 {≥}{{$\ge$}}1 {∧}{{$\wedge$}}1
    {¬}{{$\neg$}}1 {∈}{{$\in$}}1 {₁}{{$_1$}}1 {₂}{{$_2$}}1 {⊕}{{$\oplus$}}1
}

\newtheorem{theorem}{Theorem}
\newtheorem{lemma}[theorem]{Lemma}
\theoremstyle{remark}
\newtheorem{remark}[theorem]{Remark}

\begin{document}

\title[Foundations for a fast Mathlib multiplier]{Foundations for a fast integer multiplier\\ on Mathlib's Turing machines}

\author{Dan Brown$^*$}
\email{jdb1729@gmail.com}
\thanks{$^*$Author of \emph{The Da Vinci Code}.}

\date{October 11, 2026}

\begin{abstract}
We reduce the Mathlib statement \texttt{MathlibTM2.Statement} (integer multiplication in time
$O(n(\log n)^{1-\kappa})$ on \texttt{Turing.FinTM2}) to the existence of a single structured
host program that multiplies long inputs within that time. Toward that program we prove in
Lean: the complete numeric algorithm of the paper as one concrete function whose rounded
output is the product for every large input, with pairwise coprime moduli in place of primes;
the simultaneous butterfly layer as an exact computation on Gaussian integers; the round
decomposition and error bound of the synthetic transforms; the solution of the layer's cost
recurrence; and a library of host machines with exact step counts, including the network gate,
sign, leaf, packed-addition and node-body passes of the layer and the input and output stages of
the multiplier. All theorems depend only on \texttt{propext}, \texttt{Classical.choice} and
\texttt{Quot.sound}.
\end{abstract}

\maketitle

\section{The reduction}

Everything lives under \texttt{IntegerMultBounds.MathlibTM2} and builds on the Mathlib statement
and the structured stack language \texttt{Stack.lean} of the preceding contribution. A
\emph{host program} is a procedure table in that language over the Mathlib input and output
stacks and a bank of register stacks. The predicate \texttt{HostMulWithin} says that on input
$x\#y$ the program halts in the exact final configuration required by
\texttt{TM2ComputableInTime} within the given time.

\begin{theorem}[\texttt{Fast.statement\_of\_dispatch}]
If a host program satisfies \texttt{HostMulWithin} on every input of encoded length at least
$n_0$, with time $O(n(\log n)^{1-\kappa})$ for some $\kappa>0$, then
\texttt{MathlibTM2.Statement} holds.
\end{theorem}

The proof wraps the program in a dispatcher that reads up to $n_0$ input symbols into its finite
control. A shorter input is answered from a finite table; otherwise the symbols are pushed back
and the fast program runs.

\section{The numeric algorithm}

\texttt{Mult.Pipeline.pipelineCoeffs n x y} is the paper's computation with every ingredient
fixed: the parameters $b,p,d,K,T,r,\ell,t_i,s_i,\alpha,\gamma$ of \texttt{Mult.Sizes}; the CRT
reindexing $\Phi$ of \texttt{Mult.CRT}; the line maps $\widetilde A$, $\widetilde B_0$ with the
vendored weight tables; the chirp table; the computed complex convolution through twisted
synthetic transforms; pointwise products; and the final scaling and rounding.

\begin{theorem}[\texttt{Mult.Pipeline.pipeline\_correct}]
For all large $n$ and all $x,y<2^n$,
\[
 \mathrm{evalBase}\,(2^{b})\,(\texttt{pipelineCoeffs}\ n\ x\ y)=xy.
\]
\end{theorem}

The chain behind it is the vendored \texttt{ExactRecovery.exact\_product} together with new
ingredients.

\subsection*{Coprime moduli}
No prime number theorem is needed. Let $N$ be the product of the odd numbers below $4d$, take
group A moduli $Nx+1$ near $r/2$ for consecutive even $x$, and group B moduli $Ny+2$ near $r$
for consecutive odd $y$ with $y\approx 2x$. A common divisor of two moduli has no prime factor
below $4d$ and divides one of $y-2x$, $x-x'$, $y-y'$, all nonzero and at most $4d$; so the
moduli are pairwise coprime, and for $r\ge 16d^2N+8d$ they lie in the required intervals
(\texttt{Mult.Moduli.src\_good}).

\subsection*{Errors}
The source transform $\widetilde F=2^\gamma\,\widetilde B_0\,\widetilde C\,\widetilde A$ is
within $9\cdot 2^\gamma T\log_2 T$ grid units of $RF_s\Phi^{-1}$ on the half ball
(\texttt{Mult.SourceErr}), including the weighted chirp, restricted-domain compositions and the
opposite transform by conjugation. The computed complex convolution is within
$E_{\rm ring}+9$ units (\texttt{Mult.ComplexConv}), and the synthetic convolution of the twisted
records has $E_{\rm ring}=\sqrt2\,M(3\ell+1)<5T\log_2 T$ (\texttt{Mult.SynthErr}). The forward
and opposite synthetic transforms are exactly $\ell$ rounds of normalized butterflies on one
selected bit per axis followed by per-record twiddles $y^{\mp E}$ (\texttt{Mult.Rounds}).

\subsection*{A cheaper $\widetilde B_0$}
The vendored tabled $\widetilde B_0$ runs $p$ Neumann steps. With
$\|N-I\|\le 2^{-\alpha^2\theta}$ (\texttt{opNorm\_offDiagCLM\_le\_rpow}) the pipeline uses
$\alpha^2$ steps and an off-diagonal window of $\lfloor\sqrt p\rfloor/(2\alpha)+1$ terms, with the
same error conclusion (\texttt{Mult.FastB0.resampB}$_0$\texttt{NumK\_spec}): $O(s\,\alpha\sqrt p)$ term
products per line.

\section{The simultaneous layer on Gaussian integers}

The machine stores Gaussian integers modulo $2^w$. Each network gate is doubled (its
coefficients are halves of integers, \texttt{program\_gate\_coeff}) and doubles the wires it does
not write; a child returns $2^{T_c}$ times its result and the other wires are scaled by
$2^{T_c}$. With this uniform exponent every step is a ring operation, so intermediate overflow is
harmless.

\begin{theorem}[\texttt{Layer.nodeM\_correct}, \texttt{Layer.layerZ\_correct}]
The scaled node, and by induction the recursive layer \texttt{layerZ} on $W^jR_0$ rows, computes
$2^{E}$ times the kernel $C$ on every active bit.
\end{theorem}

\begin{theorem}[\texttt{Layer.topZ\_correct}]
The full layer pipeline (phase $S$, kernels on the top chunks, rows and padding, base-$m$ pieces
through \texttt{layerZ}, phase $S$, the unit $(1-i)^D2^{-\lfloor D/2\rfloor}$, truncating shift)
returns the componentwise truncation toward zero of the normalized butterflies $H_0$ on the
selected bits.
\end{theorem}

The cost recurrence of the paper is solved in \texttt{Layer.Recurrence.layer\_cost\_bound}:
with branching $s/W<m^\sigma$ (\texttt{ComplexRank25.branching\_bound}) and the interchange
exponent $\tau$ (\texttt{BinaryInterchangeRun.uniform\_bound}), the normalized cost of a layer is
$O(d^{\lambda'})$.

\section{Host machines}

Every machine below comes with an \texttt{\_exec} theorem giving its exact final register
contents and an explicit step bound.

\begin{itemize}
\item \emph{Toolkit} (\texttt{Host}): foreign calls running any vendored \texttt{Machine.Program}
with a \texttt{HoareTime} contract on host registers; streams, block routing, word automata;
index-driven routing and merging by address codes; stable radix sort; exact binary arithmetic
(addition, comparison, multiplication, division); the vendored fast
interchange of address chunks; the vendored Schönhage--Strassen multiplier
(\texttt{ssMul\_exec\_log}).
\item \emph{Layer passes} (\texttt{Layer}, \texttt{Node}): Gaussian linear automata, scaling by
$2^s$, the doubled gate on a bank of wire registers, address-dependent phases and signs, the
leaf kernels, the packed slot-XOR passes (controlled rotations, field swaps through the fast
interchange, key attachment, exact repair by extraction, radix sort and reinsertion, top-bit
fixup), unary parameter and mask utilities, and the node body
(\texttt{LayerMachine.nodeBody\_exec\_layerZ}) with abstract factor programs.
\item \emph{Multiplier stages} (\texttt{Fast}): reading the input and writing the scaled digit
records; carry propagation and the Mathlib-encoded output.
\end{itemize}

\section{Status}

The fast host program is not yet assembled. Its remaining components are the factor programs
and the recursion of the layer machine, the synthetic transform machine, the linewise
resampling machines, the table machines, the CRT reindexing machine, the pointwise and scaling
stages, and the total cost bound.

\section{Verification}

\texttt{lake build} of the library passes, \texttt{AxiomAudit.lean} passes, and
\texttt{FastPathCheck.lean} asserts the axioms of the main theorems above.

\end{document}
